% !TeX root = ../../金融经济学_熊和平_习题.tex
% 源：微信图片_20260928151417_3591_382.jpg（教材印张 293）

\begin{problem}
\problemno{13.1} 在纯交换经济基础上，有哪些资源配置的一般拓展路径？
\end{problem}
\begin{problem}
\problemno{13.2} 储蓄或储存一般可以给消费可行集带来怎样的影响？
\end{problem}
\begin{problem}
\problemno{13.3} 生产技术通常如何影响资源配置？
\end{problem}
\begin{problem}
\problemno{13.4} 简单说明纯交换经济与生产性经济的区别与联系。
\end{problem}

\begin{problem}
\problemno{13.5} 经济中有 1 个经济人，他具有对数效用函数。期末有等概率的两个状态：状态 $a$ 和状态 $b$。经济人的初始禀赋为期初 \num{200} 个消费品，但他可以通过储存将消费品跨时转移。
\begin{enumerate}[label=(\alph*)]
    \item 当利率为 $1.05\%$ 时，分别确定时间偏好率为 $1$ 和 $0.5$ 时的最优资源配置；
    \item 当利率为 $-5\%$ 时，分别确定时间偏好率为 $1$ 和 $0.5$ 时的最优资源配置；
    \item 进一步如果该经济人期末也有禀赋，分别为 $a$ 状态下 \num{100} 个消费品、$b$ 状态下 \num{50} 个消费品，分别求解 (1)(2) 中的问题。
\end{enumerate}
\end{problem}
\begin{problem}
\problemno{13.6} 经济中有 1 个经济人，他具有对数效用函数。期末有等概率的两个状态：状态 $a$ 和状态 $b$。经济人的初始禀赋为期初 \num{200} 个消费品，他可以如下生产技术进行资源跨时转移：
\[
    y_1=f_{1,\omega}(I)=
    \begin{cases}
        \sqrt{I}, & \omega=a,\\
        2\sqrt{I}, & \omega=b.
    \end{cases}
\]
\begin{enumerate}[label=(\alph*)]
    \item 求最优投资数量；
    \item 求状态价格和最优消费计划。
\end{enumerate}
\end{problem}
\begin{problem}
\problemno{13.7} 经济中有 1 个经济人，他具有对数效用函数。期末两个状态 $a$ 和 $b$ 的概率分别为 $\pi_a$ 和 $\pi_b$。经济人的初始禀赋为期初 \num{200} 个消费品，他可以如下生产技术进行资源跨时转移：
\[
    y_1=f_{1,\omega}(I)=
    \begin{cases}
        \lambda_a I, & \omega=a,\\
        \lambda_b I, & \omega=b.
    \end{cases}
\]
证券市场上存在无风险债券，其利率为 $r_f$。
\begin{enumerate}[label=(\alph*)]
    \item 给出经济人的最优决策，包括投资决策和最优消费；
    \item 求解他的最优投资数量和消费计划；
    \item 分析最优投资与利率的关系；
    \item 假设该经济人没有任何商品禀赋，回答上述 (1) 至 (3) 问。
\end{enumerate}
\end{problem}
\begin{problem}
\problemno{13.8} 在第 3 题的背景下引进经济人 2，经济人 1 只有生产技术禀赋，经济人 2 的禀赋是期初 \num{200} 个消费品：
\begin{enumerate}[label=(\alph*)]
    \item 求两个人的最优消费计划和状态价格；
    \item 求经济人 1 的最优投资。
\end{enumerate}
\end{problem}
